\honest{On Caring}{So8res}{2014}

I cannot feel the size of large numbers. A star that holds a million Earths and a star that holds a billion both just feel big; one followed by a hundred zeroes feels bigger than a billion, but not by the factor it is, the way four apples feel like twice two. This is scope insensitivity, and it matters because the things I care about are numerous. Billions live in squalor; each death is tragic however far away, and I care about every person, but I cannot scale the caring I feel for one person by a billion. My ``care-o-meter'' was calibrated for about a hundred and fifty people, and billions are suffering. So I propose that caring should work like courage: you do the right thing ``anyway,'' without the feeling. The stakes are billions suffering now and perhaps quadrillions of future people. Saving one life would feel about as good as saving the world, but behind the similar feelings there is a whole world of difference.

Then three stories about people who give for social reasons. Alice gives \$20 to students with clipboards, out of guilt, social pressure and some altruism, after they finally corner her. Bob is too busy to pour ice water over his head for a friend's challenge, so he donates \$100 instead. Christine's sorority is racing another sorority to raise money for breast cancer, and she gives a few hundred dollars herself. Their motives have to do with the social setting and only ``tangentially'' with the cause. Ask them why they do not give all their time and money to causes they think worthwhile, and they will look at you funny and think you rude; we praise those who give everything away, but privately think them crazy. Another mindset can hit you like a freight train.

Daniel is my example of the shift. After the Deepwater Horizon spill, a student with a clipboard asks him to help save oiled birds. Instead of dismissing it, he pictures a beach where a slick young bird flops to his feet and someone says three minutes of washing would save it. He decides that saving one oiled bird is worth three minutes of his time, or \$3, in some general sense and not only because it is in front of him. Expecting his feelings to misreport large numbers, he notes that thousands of birds were oiled, so he ``shuts up and multiplies,'' and finds ``with growing horror'' that he cares about oiled birds at least as much as two months of hard work or fifty thousand dollars. Then how much does he care about factory farming, hunger, poverty, war, neglected children, the future of humanity? Far more than all his money and time. \nb{Multiplying assumes that each bird is worth the same however many birds there are. The post does not state this assumption. Its only defense is that feelings misreport large numbers; it does not consider that the value of each bird might itself fall as their number grows.}

Then Daniel sees how much else needs doing and decides that he cannot spend his time or money on birds after all, not because they are not worth it, but because the opportunity cost is too high: people are dying, and civilization's future is at stake. He does not give to the wildlife charity, or to the ALS charity, or to the breast cancer charity. What changes is that if you ask him why he does not give away all his money, ``he won't look at you funny.'' He sees that his mind was lying to him about the gravity of the problems. Even the ``less important'' causes now seem worth a life, but there are too many mountains. Before, dropping everything for ALS never crossed his mind; it was not his problem. Now everything is his problem, and only the number of more urgent things stops him. Alice, Bob and Christine fail to act because they forget to see the problems; Daniel, because there are too many. In a parenthesis I hope that he ``goes on to discover'' effective altruism.

I am not preaching about how to be a good person. Many of us know we should care about distant suffering and fail to, partly because we trust our care-o-meters and think prominent altruists were gifted with extra caring. But they are people who have learned not to trust their care-o-meters, which are broken for large numbers. The lesson: you cannot feel the caring, but you can do it; stop trusting your feelings and ``switch over to manual control.'' What do you do then? I admit: ``And I don't really know yet.'' As ``a good start'' I list four organizations. One is MIRI. \nb{The author was then a research fellow at MIRI. The post does not mention this.} Part of the answer, I think, is ``a certain sort of desperate perspective'': you would give your life to the world's 100th biggest problem, ``but you can't, because there are 99 bigger problems you have to address first.'' I am not trying to guilt anyone: philanthropy is really hard, needing money and the will to throw it at distant invisible problems, and guilt is a poor long-term motivator. I would rather you join proudly.

I do not think that Gandhi, Mother Teresa and Nelson Mandela cared more than we do. The way to act is to multiply, and to trust the numbers more than the feelings. \nb{Paul Slovic, a psychologist who studies the effect the post is about, took the title of a 2007 paper from Mother Teresa: ``If I look at the mass I will never act. If I look at the one, I will.'' That is the opposite of the post's method.}

When you do the multiplication, fighting poverty and building a better future ``deserve more resources than currently exist.'' I do not show the multiplication for either. ``There is only you, and me, and everyone else who is trying anyway.'' You cannot feel the weight of the world, but sometimes you can catch a glimpse.
